\chapter{Sprint 40 --- Đặt hàng Mobile: trạng thái xử lý + kết quả gửi đơn (2026-07-26)}

\section{Bối cảnh}

Task BA nằm ở Google Sheet tab \texttt{Tasks} \textbf{row 5} (STT 4, \emph{Loại task} = UI-Issue,
\emph{Sẵn sàng cho AI} = Yes, ưu tiên Cao, phạm vi \texttt{wujia\_portal\_sale} / Mobile Portal ---
Order \& Cart): ``\emph{Bổ sung trạng thái xử lý và kết quả gửi đơn đặt hàng trên Mobile Portal}'',
kèm Figma node \texttt{4963:2}.

\textbf{Vấn đề của luồng cũ.} \texttt{POST /portal/order/submit} chạy đồng bộ rồi redirect thẳng
\texttt{/portal/purchase-history/<id>?message=order\_submitted}. Trong lúc server tạo đơn, người
dùng \textbf{không thấy phản hồi nào} --- hoặc bấm lại nút, hoặc tưởng máy treo. Khi bị chặn vì
ngoài khung giờ thì chỉ về giỏ với một banner đỏ nhỏ, \textbf{không nói khung giờ mở lại lúc nào,
ngày nào}.

\textbf{Đọc Figma qua MCP} xác nhận \texttt{4963:2} là frame duy nhất của luồng này và
\textbf{chỉ có bản mobile}, không có bản PC:

\begin{center}
\begin{tabular}{ll}
\hline
\textbf{Màn Figma} & \textbf{Nội dung} \\
\hline
02 · Đang tạo đơn & Overlay khoá thao tác trên màn giỏ: spinner + ``Đang tạo đơn\ldots'' \\
03 · Tạo đơn thành công & Trang riêng: icon check xanh, 4 dòng thông tin đơn, 2 CTA \\
04 · Backend từ chối & Trang riêng: icon đỏ, khung giờ tiếp theo + ngày, 2 CTA \\
\hline
\end{tabular}
\end{center}

Mọi số trong Figma (\texttt{SO/2026/00718}, \texttt{21/07/2026 • 09:32}, \texttt{1.352.000 đ},
\texttt{10:00 -- 16:00}, \texttt{22/07/2026}) là \textbf{minh hoạ} --- bản dựng lấy hết từ đơn thật
và config khung giờ thật.

\textbf{3 fork đã hỏi \& chốt với chủ dự án:}
\begin{itemize}
  \item \textbf{PC giữ nguyên 100\%.} Form mobile gửi thêm hidden \texttt{flow=m}; chỉ nhánh mobile
    đi vào màn kết quả mới. PC đã có màn cảnh báo riêng từ Sprint PC-1 và đã retest --- không đụng.
  \item \textbf{Màn 04 chỉ dùng cho \texttt{ORDER\_TIME\_CLOSED}.} Các lỗi submit khác
    (\texttt{CART\_EMPTY}, \texttt{CART\_QUANTITY\_INVALID}, \texttt{CART\_HAS\_INVALID\_PRODUCT},
    \texttt{CART\_IS\_PROCESSING}, \texttt{ORDER\_CREATE\_FAILED}, \texttt{STORE\_*}) giữ nguyên
    hành vi cũ: về giỏ + banner đỏ, để người dùng sửa ngay tại chỗ.
  \item \textbf{Nhãn trạng thái tái dùng, không nhân bản} --- import \texttt{\_state\_meta} của
    trang Lịch sử.
\end{itemize}

\section{Phần A --- Helper ``khung giờ tiếp theo''}

\texttt{wujia\_portal\_order\_window/models/res\_config\_settings.py} --- thêm \textbf{2 method
read-only mới}, \textbf{không sửa} rule \texttt{\_is\_within\_order\_window} hiện có:

\begin{itemize}
  \item \texttt{\_user\_now\_dt()} --- ``hôm nay'' theo timezone người dùng (cùng cách
    \texttt{\_user\_now\_hours} đang tính giờ, để ngày và giờ không lệch nguồn).
  \item \texttt{\_next\_order\_window(area\_id=None)} --- gọi lại chính
    \texttt{\_is\_within\_order\_window()} để lấy \texttt{window['windows']} (đã xử lý ưu tiên
    khu vực $\to$ fallback global), rồi chọn khung có \texttt{(day\_offset, from)} nhỏ nhất:
    \texttt{now < from} $\Rightarrow$ mở \emph{hôm nay}, ngược lại $\Rightarrow$ \emph{ngày mai}.
    Trả \texttt{\{from, to, name, date, is\_today\}} hoặc \texttt{None} khi tắt giới hạn khung giờ.
\end{itemize}

Model chỉ trả \textbf{dữ liệu thô}; việc format chuỗi để ở controller (tái dùng
\texttt{\_float\_to\_hhmm} đã có) --- giữ đúng ranh giới model/view.

\section{Phần B --- Controller}

\texttt{wujia\_portal\_sale/controllers/portal.py}:

\begin{itemize}
  \item \texttt{\_is\_mobile\_flow(post)} --- 1 dòng, dùng ở 3 điểm rẽ nhánh.
  \item \texttt{portal\_order\_submit} \textbf{chỉ đổi đích redirect}: toàn bộ logic khoá giỏ
    \texttt{FOR UPDATE NOWAIT}, validate, huỷ báo giá cũ, clear giỏ, \texttt{\_publish\_cart\_event}
    giữ nguyên từng dòng. Thành công: mobile $\to$ \texttt{/portal/order/submitted/<id>};
    PC $\to$ URL cũ. \texttt{ORDER\_TIME\_CLOSED} gom về \texttt{\_submit\_time\_closed(post)}
    dùng chung cho \textbf{cả 2 nhánh} (check trước khi create \textbf{và}
    \texttt{ValidationError} chứa ``khung giờ'' từ tầng \texttt{sale.order.create}).
  \item \textbf{\texttt{GET /portal/order/submitted/<int:order\_id>}} --- guard scope: đơn phải tồn
    tại, \texttt{is\_portal\_order}, \texttt{franchise\_id} khớp cửa hàng đang chọn; sai $\to$ đẩy về
    \texttt{/portal/order} (không tiết lộ đơn có tồn tại hay không). Chặn thêm
    \texttt{state == 'cancel'}: \texttt{SALE\_STATE\_META} không map \texttt{cancel} nên sẽ rơi về
    nhãn an toàn ``Đang xử lý'' --- hiển thị đơn đã huỷ như một ``kết quả gửi đơn'' là sai nghiệp vụ.
    \texttt{date\_order} lưu UTC $\to$ quy về tz người dùng bằng
    \texttt{fields.Datetime.context\_timestamp} để ``09:32'' đúng giờ cửa hàng.
  \item \textbf{\texttt{GET /portal/order/rejected}} --- whitelist \texttt{reason} chỉ nhận
    \texttt{ORDER\_TIME\_CLOSED}; mã khác đẩy về giỏ kèm banner, mã lạ về
    \texttt{internal\_error} --- không reflect chuỗi tuỳ ý từ query string ra UI, đúng nguyên tắc
    của \texttt{\_resolve\_messages}.
  \item \textbf{Dependency mới:} \texttt{wujia\_portal\_purchase\_history} vào \texttt{depends} của
    \texttt{wujia\_portal\_sale} để import \texttt{\_state\_meta}
    (\texttt{draft} $\to$ \emph{Chờ xác nhận}/\texttt{pending}). Không có vòng phụ thuộc ---
    \texttt{purchase\_history} chỉ depend \texttt{wujia\_sale}, \texttt{wujia\_portal\_base},
    \texttt{wujia\_delivery} --- và CTA ``Xem chi tiết đơn hàng'' trỏ thẳng sang trang đó nên phụ
    thuộc này đúng ngữ nghĩa. Màn kết quả và trang lịch sử vì thế \textbf{luôn nói giống nhau}.
\end{itemize}

\section{Phần C --- Template, CSS, JS}

\textbf{\texttt{views/portal\_order\_result.xml} (mới, 124 dòng).} Một khung dùng chung
\texttt{mres\_shell} (icon + tiêu đề + phần thân + CTA) cho cả 2 màn, đặt trong
\texttt{wujia\_portal\_layout.app\_layout} + \texttt{wj\_page\_header}
(\texttt{ph\_platform='m'}, \texttt{ph\_variant='title'}) $\to$ header, Current Store Strip và
bottom nav lấy từ shell sẵn có, không dựng lại. Box thông tin dùng \texttt{dl/dt/dd} thay vì
\texttt{label} --- tránh rule \texttt{label \{padding-left:.2rem\}} của \texttt{bootstrap-extended}
(§L4).

\textbf{\texttt{views/portal\_order\_cart.xml}.} Thêm hidden \texttt{flow=m} vào form mobile và
overlay \texttt{\#wj-order-submitting}. Overlay đặt \textbf{ngoài} \texttt{.wujia-mcart}: cart-sync
swap \texttt{innerHTML} của panel đó mỗi lần reconcile, để bên trong sẽ bị xoá mất giữa chừng.

\textbf{CSS (+242 dòng \texttt{\_components.css}, +10 token \texttt{\_variables.css}).} Hai prefix
\texttt{.wujia-msubmit-*} và \texttt{.wujia-mres-*}, mọi rule scope trong 2 prefix đó. Token
append-only, không sửa token cũ $\to$ blast radius = 0. Theo §L4, tiêu đề màn kết quả là
\texttt{<h2>} nên bị rule tag-level \texttt{h1, h2 \{font-size: … !important\}} của shell nuốt:
trung hoà bằng \texttt{.wujia-mres .wujia-mres-title} + \texttt{!important}
\textbf{trong phạm vi component}, không đụng file shared. Spinner tôn trọng
\texttt{prefers-reduced-motion}. CSS nằm ở \texttt{wujia\_portal\_layout} (nạp qua \texttt{<link>}
thủ công) nên \textbf{phải bump} \texttt{?v=1155 $\to$ 1157} --- §9 gotcha \#1.

\textbf{JS (+48 dòng \texttt{portal\_order.js}).} Bắt \texttt{submit} trên form mobile (nhận diện
bằng \texttt{input[name=flow][value=m]}, delegation trên \texttt{document} vì panel giỏ bị swap):
lần 1 set cờ + \texttt{disabled} nút + bật overlay, lần 2 \texttt{preventDefault()}. Nút submit
không có \texttt{name} nên disable không làm mất payload; \texttt{portal\_note} vẫn gửi vì nằm cùng
form (FUNC-MOB-ORDER-005). \texttt{pageshow} + \texttt{ev.persisted} $\to$ tắt overlay, xoá cờ, mở
lại \textbf{đúng nút mà chính JS đã khoá} --- nút bị server khoá vì ngoài khung giờ (WJ-ORD-006)
phải giữ nguyên trạng thái khoá.

Lớp chống trùng \textbf{thật} vẫn ở server (\texttt{FOR UPDATE NOWAIT} $\to$
\texttt{CART\_IS\_PROCESSING}; lần 2 sau khi giỏ đã clear $\to$ \texttt{CART\_EMPTY}); JS chỉ là
lớp UI.

\section{Gì đã làm (nghiệp vụ)}

Cửa hàng bấm ``Gửi đơn đặt hàng'' trên điện thoại giờ \textbf{thấy hệ thống đang làm việc}: màn hình
mờ lại kèm vòng quay và dòng ``Đang tạo đơn\ldots{} Vui lòng không đóng màn hình'', nút gửi khoá lại
nên không thể bấm nhầm hai lần thành hai đơn.

Gửi thành công, họ vào thẳng một trang xác nhận thay vì bị đẩy sang danh sách lịch sử: dấu tích xanh,
\textbf{mã đơn, thời gian đặt theo giờ cửa hàng, tổng tiền và trạng thái} --- trạng thái dùng đúng
nhãn của trang Lịch sử nên hai chỗ không bao giờ nói khác nhau --- kèm 2 lối đi tiếp: xem chi tiết
đơn vừa tạo, hoặc quay lại chọn thêm sản phẩm.

Gửi lúc ngoài khung giờ, thay vì một dòng banner đỏ mơ hồ, họ nhận một trang nói rõ ba điều: đơn
\textbf{chưa} được tạo, \textbf{giỏ hàng vẫn còn nguyên}, và \textbf{khung giờ đặt hàng tiếp theo mở
lúc mấy giờ, ngày nào} --- lấy theo khung giờ của chính khu vực cửa hàng đó, có ghi rõ ``Hôm nay''
khi khung mở lại trong ngày. Nhờ vậy cửa hàng biết chờ tới lúc nào thay vì bấm thử lại nhiều lần.

\section{Lý do/Trade-off}

\begin{itemize}
  \item \textbf{Tách nhánh bằng \texttt{flow=m} thay vì detect user-agent.} Một trường hidden trong
    form là thứ duy nhất khác nhau giữa 2 luồng; PC không gửi $\to$ hành vi PC bất biến theo đúng
    nghĩa đen, không phụ thuộc heuristic thiết bị.
  \item \textbf{Màn 04 chỉ cho ngoài khung giờ.} Các lỗi còn lại người dùng \emph{sửa được ngay tại
    giỏ} (số lượng sai, sản phẩm ngừng bán\ldots) --- đẩy sang trang riêng rồi bắt bấm quay lại là
    thêm một bước vô ích.
  \item \textbf{Không tạo màn PC.} Figma \texttt{4963:2} không có bản PC; dựng thêm sẽ là suy diễn
    thiết kế, trái §7 ``ask-don't-assume''.
  \item \textbf{Card kết quả có \texttt{max-width} 420px.} Figma chỉ vẽ 359px trên khung 391;
    giới hạn để khi mở 2 URL này trên màn rộng (chúng dùng chung route, không chặn theo viewport)
    card không giãn hết bề ngang trông vỡ bố cục.
  \item \textbf{Không migration.} Không thêm field nào $\to$ chỉ bump version 3 module.
\end{itemize}

\section{Verify}

\texttt{-u wujia\_portal\_sale,wujia\_portal\_order\_window,wujia\_portal\_layout --stop-after-init}:
\textbf{RC=0}, \texttt{Modules loaded} + \texttt{Registry loaded}, \textbf{0 ERROR/CRITICAL/Traceback}
trong \texttt{logs/odoo.log}.

Đo computed-style bằng Playwright/chromium ở \textbf{391$\times$844} (không nhìn ảnh, theo §L4):

\begin{center}
\begin{tabular}{lll}
\hline
\textbf{Thành phần} & \textbf{Đo được} & \textbf{Ghi chú} \\
\hline
Tiêu đề kết quả & 21px / 700 / \texttt{\#111827} & thắng được rule \texttt{h2 !important} \\
Card kết quả & 357.4px & Figma 359 trên khung 391 \\
Vòng icon ngoài / trong & 80 / 58px & \\
CTA chính / phụ & cao 48 / 44px & \\
Box 4 dòng thông tin & 323.4$\times$173, hàng 40px & \\
Box khung giờ (đỏ) & 323.4$\times$109.5 & \\
Giá trị khung giờ & 19px / 700 / \texttt{\#DC2626} & \\
Badge trạng thái & 13px, nền \texttt{\#EAF8EF} chữ \texttt{\#16A34A} & đơn \texttt{sale} = Đã xác nhận \\
Card overlay & 313$\times$178, tiêu đề 20px/700 & \\
\hline
\end{tabular}
\end{center}

Guard scope kiểm bằng URL trực tiếp: chưa chọn cửa hàng $\to$
\texttt{/portal/order?error=STORE\_NOT\_SELECTED}; chọn đúng cửa hàng $\to$ render đủ dữ liệu thật
của đơn trong DB. \texttt{/portal/order/rejected?reason=ORDER\_TIME\_CLOSED} render đúng box khung
giờ tiếp theo.

\textbf{Chưa deploy UAT.} Lệnh cần chạy trên server:
\texttt{git pull} $\to$
\texttt{python odoo-bin -c <conf> -d wujia\_tea\_19 -u wujia\_portal\_sale,wujia\_portal\_order\_window,wujia\_portal\_layout --stop-after-init}
$\to$ restart service. Không có module mới nên \textbf{không cần} \texttt{-i} (§6 gotcha).
